\centerline{\bf \Huge Суммы цифр}

\begin{flushright}
        \scriptsize{Эпизод 14. Seele, трон для души}
\end{flushright}

\textit{Пусть $S(n)$ - сумма цифр натурального числа $n$.}

\problem{1}
Докажите, что $S(a+b)=S(a)+S(b)-9 p$, где $p-$ количество переносов при сложении $a$ и $b$ в столбик.

\problem{2}
Найдите наименьшее натуральное $n$, для которого $S(n)=13$, а также
(a) $S(4 n)=52$;
(б) $S(4 n)=34$.

\problem{3}
Известно, что сумма цифр натурального числа $n$ равна 100 , а сумма цифр числа $5 n$ равна 50. Докажите, что $n$ четно.

\problem{4}
Докажите, что

$$
S(1)-S(2)+S(3)-\ldots-S\left(7^{100}-1\right)+S\left(7^{100}\right)=\frac{7^{100}+1}{2} .
$$


\problem{5}
Пусть $p(n)$ - количество переносов при сложении чисел 2022 и $2022 \cdot n$ в столбик. Например, $p(1)=p(2)=p(3)=0, p(4)=3$. Докажите, что

$$
p(1)+p(2)+\cdots+p\left(10^{2022}-1\right)=\frac{2}{3} \cdot\left(10^{2022}-1\right) .
$$


\problem{6}
Антиподом натурального числа называется число, записанное теми же цифрами в обратном порядке. Каждое $k$-значное натуральное число без нулей в записи сложили с его антиподом и выписали на доску все результаты. Оказалось, что на доске нет чисел, в записи которых были бы только нечетные цифры. Какие остатки может давать $k$ при делении на 4 ?

\problem{7}
Про натуральное число $n$ известно, что $S(n)=100$, а $S(55 n)=1000$. Найдите $S(3 n)$.

\problem{8}
Найдите наименьшее натуральное $n$, для которого $S(n)$ и $S(n+10)$ делятся на 19 .

\problem{9}
Федя задумал натуральное число $n$ и сообщил Эрдему его сумму цифр $S(n)$. Эрдем хочет угадать число Феди. Он может выбирать любое натуральное число $m$ и узнавать у Феди сумму цифр числа $|n-m|$. Как Эрдему за $S(n)$ вопросов гарантированно узнать число Феди?